% Section 2: Program processes and the arity-one carrier.

\subsection{The native object}

\begin{definition}[Open finite program]\label{def:open-program}
An \emph{open finite program} $P$ consists of
\begin{enumerate}
\item a finite typed expression DAG $B(P)$, its \emph{body}, whose nodes are
  operations and whose edges carry typed wires;
\item a finite set $\partial P$ of \emph{holes}, each a named position of declared
  type, which are the only open leaves of $B(P)$;
\item a labelling that keeps distinct, for every source in $B(P)$, its
  \emph{occurrences}: sharing and discarding are explicit and are not quotiented
  away;
\item the construction history of $B(P)$: the nesting of scopes and the partial
  order in which nodes were introduced.
\end{enumerate}
\end{definition}

Composition grafts an argument program into a hole; the result is again an open
finite program.  Evaluation --- numerical, symbolic, spectral --- is an observation
of such a program under a declared policy, not its identity.  Consequently the
implications
\[
\text{same value}\Rightarrow\text{same program},\qquad
\text{same differential}\Rightarrow\text{same program}
\]
both fail, and the failure is recorded rather than repaired by quotienting.

\begin{definition}[Construction order]\label{def:construction-order}
The \emph{construction order} of $P$ is the partial order $\preceq_P$ on the
operation nodes of $B(P)$ in which $u\preceq_P v$ when $u$ must be constructed
before $v$: dependency edges, scope nesting, and source/occurrence creation.
\end{definition}

\begin{definition}[Evaluation schedule]\label{def:schedule}
An \emph{evaluation schedule} of $P$ is a linear extension of $\preceq_P$, that is,
a total order $\sigma$ on the operation nodes such that $u\preceq_P v$ implies
$u$ precedes $v$ in $\sigma$.  The set of schedules is written $\Sch(P)$.
\end{definition}

Two schedules of one program may differ while the program does not.  A schedule is
therefore extra data; it is the process analogue of a chosen linear extension of a
dependency poset, and it is not the construction order itself.

\subsection{Three readings of one process}

One program process is read in three coordinated ways.

\begin{itemize}
\item \textbf{Construction $K$}: which open programs are grafted into which holes,
  the nesting of scopes, explicit copy and discard, and the dependency partial
  order generated by those events.
\item \textbf{Causal development $C$}: the orientation of construction from a
  completed past toward enabled future events; crossing an enabled node replaces
  the frontier wires it consumes by the wires it produces.
\item \textbf{Frontier $F$}: a downward-closed set of completed events, read as the
  checked wires crossing from that past to its future --- an occurrence-decorated
  open boundary, not a tuple of values.
\end{itemize}

The three readings are coupled by
\[
\text{construction dependency}\;\longrightarrow\;\text{legal causal advance}
\;\longrightarrow\;\text{new open cut},
\]
and it is exactly this coupling that a schedule linearizes.  The statement that $K$,
$C$ and $F$ are readings of one process, rather than three structures, is the
primitive of the theory; it is imported from the process corpus, where the finite
slices and their composition laws are checked \cite{ProcessGeometryCore2026}.

\begin{remark}[Slices]\label{rem:slices}
For downward-closed $U\subseteq V$, the slice of $P$ between the cuts $U$ and $V$
retains the events of $V\setminus U$, both cuts, the changed boundary wires, the
unchanged through wires, the occurrences and the history.  Adjacent slices compose
after input revalidation, and two distinct linear schedules of one diamond yield the
same outer slice.  The last statement is the finite form of the separation between
schedule and process: schedule path is extra data, not slice event order.
\end{remark}

\subsection{The arity-one principle}

The link to arithmetic expression geometry is the carrier of Paper~V.  Its complex
unit $\Omega=TSR$ acts on the six-state carrier
$W=\Lambda^{1}E\oplus\Lambda^{2}E$; the three planes
$W_T=\langle T,SR\rangle$, $W_S=\langle S,RT\rangle$, $W_R=\langle R,TS\rangle$ pair
each reading with the oriented relation of the other two, and the chirality sign law
states that cyclic permutations preserve $\Omega$ while transpositions reverse it
\cite{Yuan2026AEGPaperFive}.

\begin{proposal}[Arity one]\label{prop:arity-one}
The iota carrier acts at one slot.  Reading a program's holes as the slots at which
the carrier acts, the carrier induces on the holes the order in which they are
consumed, and through that order an order on the operations that fill them.  The
spine of a sequential arithmetic expression, and the left and right combs that
decompose a general expression tree, are then the single-spine and path cases of one
slot order rather than independent primitives.
\end{proposal}

The proposal is structural: it names the mechanism by which the carrier reaches the
combinatorics, and it is the organising principle of this paper.  What it does not by
itself assert is that every process order is realized by the carrier; that is the
content of the equivalence statements of \S\ref{sec:equivalence}, each with its own
status.

\subsection{The sequential case}

\begin{definition}[Spine and combs]\label{def:spine-combs}
A finite binary expression tree is \emph{sequential} when its internal-node
dependency poset has a unique linear extension; equivalently, when its internal
vertices form a single spine \cite{Yuan2026AEGFoundations}.  In the marked spinal
history of Paper~I, the accumulator evolves along that spine and each step carries a
one-hole context $C_{\omega,c}^{(1)}[z]=\omega(z,c)$ or
$C_{\omega,c}^{(2)}[z]=\omega(c,z)$; the sequence of slot choices is the chirality
word.  A \emph{left comb} is the spine whose contexts are all left slots, and a
\emph{right comb} the spine whose contexts are all right slots.
\end{definition}

\begin{proposition}[Sequential histories are single-spine programs]\label{prop:sequential-leaf}
Every marked spinal history of Paper~I is an open finite program with exactly one
hole and a construction order whose unique linear extension is the history itself.
Conversely, an open finite program with one hole whose body is a chain of unary
one-hole contexts, up to the slot convention of its contexts, is a marked spinal
history.
\end{proposition}

\begin{proof}
Immediate from the definitions: the body of the history is the chain of contexts,
the single hole is the accumulator position, the construction order is the chain
order, and the unique linear extension is the chronological order.  The slot
convention is exactly the chirality word.  The converse reads the same data in the
other direction; the two constructions are inverse up to the declared convention on
slot labels.
\end{proof}

\begin{remark}[What the sequential case does not see]\label{rem:sequential-blind}
A sequential history has one schedule, so the distinction between a process and its
linearization is invisible there.  Everything this paper adds begins with programs
that have more than one schedule: branching bodies, multiple holes, and the
interleavings that a general dependency order admits.  In this sense Paper~I is the
single-leaf theory and the present paper is the tree-level one.
\end{remark}

\subsection{The poset layer}

Fixing the sequential leaf does not yet fix the order theory of the branching case.
The dependency poset of a general program is a tree-shaped poset, and its linear
extensions are counted by the tree's hook-length formula rather than by a single
chain.  Two independent structures must be kept apart:

\begin{itemize}
\item the \emph{planar} structure, which records left and right slots and is carried
  by the corner order of \S\ref{sec:orders};
\item the \emph{poset} structure, which records only dependency and forgets
  left/right, and is carried by the down-set lattice $O(P)$ of Paper~I.
\end{itemize}

The distinction is not academic: the Donaghey-type transforms act on the planar
structure and do not descend to the poset layer, a point developed in
\S\ref{sec:orders} and registered as an open direction in
\S\ref{sec:boundaries}.  Keeping the two layers apart is a standing requirement of
this paper.
